% Copyright (c) 2026 Eugene Joseph M. Ragasa.
% SPDX-License-Identifier: Apache-2.0
% Author: Eugene Joseph M. Ragasa
% Formal proof completed: 2026-10-19
\documentclass[11pt]{article}
\usepackage[T1]{fontenc}
\usepackage{amsmath,amssymb,booktabs}
\usepackage[margin=1in]{geometry}
\usepackage[hidelinks]{hyperref}
\usepackage{tabularx}
\title{Proof Dossier: Two Unitary Alignments Have Common Power}
\author{Eugene Joseph M. Ragasa}
\date{Formal proof completed: 2026-10-19}
\begin{document}
\maketitle
\section{Identity and status}
\begin{tabularx}{\textwidth}{@{}l>{\raggedright\arraybackslash}X@{}}
\toprule
Statement ID & \texttt{BG-BLOCH-THM-007} \\
Status & \texttt{proved} \\
Lean declaration & \path{ProjectKoios.BandGauge.Bloch.unitaryConjugateSampledFields_commonPower} \\
Module & \path{ProjectKoios/BandGauge/Bloch/PeriodicRankTwo.lean} \\
\bottomrule
\end{tabularx}

\section{Exact statement}
If the sampled centered parent field has weighted power $P(D)=\kappa$, then any two
pointwise unitary fields $U_i,V_i$ give
\[
 P(U^\dagger DU)=\kappa,
 \qquad P(V^\dagger DV)=\kappa.
\]

\section{Source and proof}
This source-normalized bridge \texttt{adapts} standard unitary invariance to the common-power premise for $D_{i,\phi}$ and $B_i$ in
lines 603--624 of the Ragasa2026A v3 manuscript. Owner: \path{ksdft2effmass}; immutable
commit
\begin{center}\small\texttt{99a46cb59d0c6f27ada7d709f194996ef9619d9f}.\end{center}
Apply BG-BLOCH-THM-006 independently to the two unitary fields and substitute the parent
power equality.

\section{Boundary and trusted base}
The underlying unitary-invariance fact carries no novelty claim. The theorem does not identify $U$ and $V$ with the publication's concrete alignments.
Lean and pinned Mathlib report \texttt{propext}, \texttt{Classical.choice}, and
\texttt{Quot.sound}. No admission or scientific axiom is used. The proof and prose were
AI-assisted and kernel-checked. No independent expert review is recorded. Status:
\texttt{proved}; working draft.
\end{document}
